\documentclass[11pt]{scrartcl}

%---------------------------------
% Sprache, Schriften, Zeichensatz
%---------------------------------
\usepackage[ngerman]{babel}

\usepackage[T1]{fontenc}
\usepackage[utf8]{inputenc}

\usepackage{csquotes}	% für babel

%---------------------------------
% Datum/Zeit
%---------------------------------
\usepackage[ngerman]{datetime}

\newdateformat{germandate}{\THEDAY. \monthname[\THEMONTH] \THEYEAR}


%---------------------------------
% BibLaTeX: Online-Quellen
%---------------------------------
\usepackage[backend=biber, style=numeric, sorting=none]{biblatex}
% sorting=none -> keine Sortierung, Standard ist alphabetisch
\addbibresource{literatur.bib}
% TexMaker-Kommando für bib(la)tex: "biber" %
% (Original-Einstellung für bibtex: "bibtex" %.aux
%---------------------------------
% Meta Variables
%---------------------------------
\newcommand{\MetaInstitute}{Hochschule Bremen}
\newcommand{\MetaUnit}{Fakultät 4 -- Elektrotechnik und Informatik}
\newcommand{\MetaTitle}{Bachelorthesis}
\newcommand{\MetaSubtitle}{Transformation einer Bestandsanwendung in eine Cloud-Native-Architektur
}
\newcommand{\MetaTask}{Exposé}
\newcommand{\MetaAuthorName}{Lukas}
\newcommand{\MetaAuthorSurname}{Karsten}
\newcommand{\MetaAuthor}{\MetaAuthorName~\MetaAuthorSurname}
\newcommand{\MetaStudentNumber}{\textit{5011712}}
\newcommand{\MetaStudyProgram}{Internationaler Studiengang Medieninformatik (B.Sc.)}
% \newcommand{\MetaCoAuthor}{mit Helmut Eirund}
%\newcommand{\MetaDate}{\germandate\today}
\newcommand{\MetaDate}{04.\ November 2021}
\newcommand{\MetaVersion}{1.01}


%---------------------------------
% Querverweise
%---------------------------------
% Hyperref (sollte unbedingt vor geometry-Paket geladen werden)
% für Anzeige in Acrobat (pdfstartview)
\usepackage[bookmarks=false,
    pdfstartview=FitV,
    pdfhighlight={/I},
	colorlinks = true,
	linkcolor = blue,
	urlcolor  = blue,
	citecolor = blue,
    pdfborder={0 0 0},
    german]{hyperref}

\hypersetup{
  pdfauthor   = {\MetaAuthor},
  pdftitle    = {\MetaTitle},
  % pdfsubject  = {\MetaUnit, \MetaTask},
  pdfsubject  = {\MetaTask},
  pdfkeywords = {\MetaTitle, \MetaUnit, \MetaInstitute},
  % pdfkeywords = {\MetaTitle, \MetaUnit, \MetaTask, \MetaInstitute},
}

% Für Verweise mit Angabe des Typs des referenzierten Objekts (z.B. "Kapitel")
\usepackage[ngerman]{cleveref}

\crefname{section}{Kapitel}{Kapitel}	% Anpassung der Typbezeichnungen
\crefname{subsection}{Abschnitt}{Abschnitte}
\crefname{lstlisting}{Listing}{Listing} % neue Typbezeichnung


%---------------------------------
% Grafiken, Farben
%---------------------------------
\usepackage{graphicx}
\graphicspath{ {figures/} }

\usepackage{xcolor}

\definecolor{keyword}{HTML}{0000FF}	% Farben für Listings
\definecolor{string}{HTML}{D12F1B}
\definecolor{comment}{HTML}{008400}
\definecolor{lightgrey}{rgb}{0.99,0.99,0.99}

\definecolor{note}{rgb}{1,1,0.8}		% Farbe für Notizen


%---------------------------------
% Seitenlayout: Abmessungen
%---------------------------------
% Option für zusätzlichen Rand zum Binden: bindingoffset
% Default-Verhältnis von oberem (innerem) zu unterem (äußerem) Rand: 2:3
%\usepackage[top=3.6cm,bottom=4.50cm,left=2.3cm,bindingoffset=0.5cm, pdftex,twoside,a4paper]{geometry}
\usepackage[top=3.6cm,bottom=4.50cm,left=2.3cm,bindingoffset=0.5cm, pdftex,a4paper]{geometry}

\setlength{\parindent}{6mm} 
\setlength{\parskip}{0.2cm} 


% Raum für Notizen mit note-Paket
%\usepackage[top=3.6cm,bottom=4.5cm,right=4cm,bindingoffset=0.5cm,twoside]{geometry}

% "includemp" zieht Notizbereich (marginpar) von Druckbereich ab:
%\usepackage[top=3.6cm,bottom=4.50cm,left=2.3cm,bindingoffset=0.5cm,marginparwidth=3cm, includemp, pdftex,twoside,a4paper]{geometry}

%---------------------------------
% Seitenlayout: Kopf- und Fußzeilen
%---------------------------------
\usepackage[headsepline]{scrlayer-scrpage}

\clearpairofpagestyles
\automark[section]{section}	% Anzeige der "section" in der Kopfzeile
\lohead{\headmark}
\cofoot[\pagemark]{\pagemark} % Anzeige der Seitenzahl in der Fußzeile
\pagestyle{scrheadings}

%\pagestyle{empty}


%---------------------------------
% Seitenlayout: Überschriften
%---------------------------------
\addtokomafont{subsubsection}{\normalfont\sffamily}	% subsubsection nicht bold
\addtokomafont{paragraph}{\normalfont\sffamily}	% paragraph nicht bold


%---------------------------------
% Enumeration mit Zahlen auf allen Ebenen (für Gliederung)
%---------------------------------
\usepackage{enumitem}

\newlist{gliederung}{enumerate}{4}
\setlist[gliederung]{label*=\arabic*.}


%---------------------------------
% Listings
%---------------------------------
\usepackage{listings}

\lstset{
  language=Java,
  basicstyle=\ttfamily,
  showstringspaces=false, % lets spaces in strings appear as real spaces
  columns=fixed,
  keepspaces=true,
  keywordstyle=\color{keyword},
  stringstyle=\color{string},
  commentstyle=\color{comment},
  frame=tb,	% Rahmen = single
%  framerule=1pt,
  showstringspaces=false,
  basicstyle=\footnotesize\ttfamily,
  backgroundcolor=\color{lightgrey},
  numbers=left
}


%---------------------------------
% To Do Notes
%---------------------------------
\usepackage[textwidth=3.5cm, backgroundcolor=note]{todonotes}


%---------------------------------
% Tabellen
%---------------------------------
\usepackage{multirow}
%usepackage{array}
\usepackage{makecell}
\usepackage{booktabs}	% http://ctan.org/pkg/booktabs

\newcommand{\tabitem}{~~\llap{\textbullet}~~}


%---------------------------------
% Floating environments genauer positionieren
%---------------------------------
\usepackage{float}

%---------------------------------
% Wasserzeichen
%---------------------------------
%\usepackage{draftwatermark}
%
%\SetWatermarkText{ENTWURF}
%\SetWatermarkScale{2.5}


%---------------------------------
% Commands
%---------------------------------
\newcommand{\HRule}{\rule{\linewidth}{0.2mm}}	% Horizontal line for title page
\newcommand{\qto}[1]{\glqq #1\grqq}				% Anführungszeichen

\newcommand{\urlMitUmlauten}[2]{\texttt{\href{#1}{#2}}}				% URL mit Umlauten

% Abkürzungen (mit korrekten Abständen)
\usepackage{xspace}
\newcommand{\zB}{\mbox{z.\,B.}\xspace}
\newcommand{\dH}{\mbox{d.\,h.}\xspace} % Kommando \dh ist schon definiert
\newcommand{\ggf}{ggf.\xspace}
\newcommand{\evtl}{evtl.\xspace}
\newcommand{\bzw}{bzw.\xspace}


%---------------------------------
% Document start
%---------------------------------
\begin{document}

%---------------------------------
% Titlepage
%---------------------------------
\begin{titlepage}
  	\shortdate % Use Short Date
  	\center % Center everything on the page

  	~\\[1cm]

	%---------------------------------
	% HEADER SECTIONS
	%---------------------------------

	\begin{figure}[h!]
    		\centering
    		\resizebox{10cm}{!}{
      		\includegraphics[width=90mm,keepaspectratio]{HSB_Horizontal_RGB}
    		}
	\end{figure}

	\vspace{-0.5cm}
	\textsc{\Large \MetaInstitute}\\[0.2cm] % Major heading such as course name
	\textsc{\Large \MetaUnit}%[2.5cm] % Major heading such as course name
	
	\textsc{\large \MetaStudyProgram}\\[1.5cm]
	
	%---------------------------------
	% DOCUMENT TYPE SECTION
	%---------------------------------
	\textsc{\LARGE \MetaTask}\\[1.5cm] % Minor heading such as course title

	%---------------------------------
	% TITLE SECTION
	%---------------------------------
%	\HRule \\[0.5cm]
%	{
%		\LARGE \bfseries \MetaTitle \\[0.50cm] % Title of your document
%		\par
%	}
%	\HRule \\[1.5cm]
	\HRule \\[0.5cm]
	{
		\LARGE \bfseries \MetaTitle \\[0.50cm] % Title of your document
		\Large \bfseries -- \MetaSubtitle\ -- \\[0.50cm] % Title of your document
		\par
	}
	\HRule \\[1.5cm]

	%---------------------------------
	% AUTHOR SECTION
	%---------------------------------
	\large 
	\MetaAuthor\ (\MetaStudentNumber)\
	% \MetaCoAuthor\\[0.25cm]

	%---------------------------------
	% DATE SECTION
	%---------------------------------
	\vspace*{\fill}
	{
     \large \MetaDate\ (Version \MetaVersion)
	}
\end{titlepage}




%---------------------------------
% EINLEITUNG
%---------------------------------
\section{Einleitung}

Moderne IT-Systeme folgen bei ihrer Bereitstellung % Bei ihrer Bereitstellung, also ihrem Deployment? Nicht eher bei ihrem Design oder ihrem gesamtne Lifecycle?
oft den Prinzipien sogenannter Cloud-Native-Architektur und sind damit so gestaltet, dass sie skalierbar, resilient, administrierbar, observierbar und automatisiert sind \cite{Lee:2021}. 
Allerdings kommt es selten vor, dass diese Systeme von Grund auf neu entwickelt werden. Ein Großteil der essentiellen Geschäftsprozesse bestehender Unternehmen läuft bereits in Systemen, bei denen diese architektonischen Erkenntnisse noch keine Anwendung fanden. Eine große Herausforderung ist es daher heute, die bestehende Infrastruktur schrittweise so zu verändern, dass möglichst viele der modernen Prinzipien angewendet werden können, ohne den laufenden Betrieb zu stören oder eine komplette Neuentwicklung der Server-Architektur zu erzwingen \cite{Fiedelholtz:2021}. %Ich würde hier noch klarer raus stellen, dass die Software nicht neu geschrieben werden kann, wegen der besetehenden Business Logik und man nun quasi versucht es in der Infrastruktur zu lösen.
Vor dieser Aufgabe steht auch das Unternehmen \textit{Startnext}, eine Online-Plattform zum Crowdfunding von kreativen Ideen, nachhaltigen Projekten und Start-Ups. Im Rahmen dieser Arbeit wird ein Konzept erarbeitet, um die Startnext-Plattform von der bestehenden, klassischen Architektur schrittweise in eine Cloud-Native-Infrastruktur zu überführen und die Vorteile dieser Technologie %Architektur? oder moderne Technologien?
effektiv zu nutzen. Die erarbeiteten Methoden werden praktisch umgesetzt, diskutiert und der Erfolg der Maßnahmen bewertet.

%---------------------------------
% PROBLEMSTELLUNG UND LÖSUNGSANSATZ
%---------------------------------
\section{\label{sec:problem_loesung}Problemstellung und Lösungsansatz}
%%---------------------------------
%% Problemstellung
%%---------------------------------
\subsection{Problemstellung}
Moderne Softwareentwicklung stellt hohe Anforderungen an ihre Infrastruktur. Um neue Features zu testen, werden kurzzeitig Testinstanzen benötigt, Änderungen sollen mehrmals täglich ins Produktivsystem eingespielt werden \cite{IBM:2019}. Lastspitzen müssen schnell abgefangen und auf schwerwiegende Bugs muss mit dem Zurückrollen zu vorherigen Versionen reagiert werden \cite{Mikkelsen:2019}. Eine klassische Bereitstellung von Software über selbst- oder fremdgehostete Server schränkt dabei die Anpassungsmöglichkeiten stark ein und lässt Änderungen an der Infrastruktur nur sehr träge zu, wenn zum Beispiel ein Bedarf an mehr Rechenleistung, weiteren Rechen-Instanzen oder Serverstandorten besteht. %man könnte hier auch noch die Probleme der Kapazitätsplanung Kosten vs. Verfügbarkeits Tradeoff anbringen mit klassicher Infrastruktur 
\newline 
Die Wartung und Aktualisierung der Infrastruktur erfolgt über das Einwählen und manuelle Aufspielen von Updates. Dabei verändert sich der Zustand jeder einzelnen Instanz über die Zeit. Es kommt zum Auseinanderdriften von Zuständen der einzelnen Server. Zusätzlich ist der Weg zu diesem Zustand schwer bis nicht nachvollziehbar. Das Updaten der Instanzen wird unvorhersehbar und potentielle Fehler passieren auf produktiven Systemen. %Mangelnde Observierbarkeit bei verteilten Systemen?

%Die Startnext Infrastruktur ist wie folgt aufgebaut und hat folgende Probleme
%%---------------------------------
%% LÖSUNGSANSATZ
%%---------------------------------
\subsection{Lösungsansatz}
Einen Lösungsansatz für genau diese Probleme bietet das Konzept von Infrastructure as a Service (IaaS). Als IaaS werden Cloud-Computing-Services bezeichnet, die Rechner-, Speicher- und Netzwerkressourcen auf Anfrage bereitstellen. Dabei wird jede der Ressourcen als separate Komponente angeboten und nur für die Dauer der Nutzung in Rechnung gestellt. Der Aufwand und die Kosten für das Betreiben von eigenen physikalischen Servern und der zugehörigen Infrastruktur werden also umgangen, während das Einrichten sowie Deinstallieren neuer Komponenten vollautomatisiert und innerhalb weniger Sekunden erfolgen kann.\newline
Durch die alleinige Umstellung zu einem Hosting bei einem Cloud-Anbieter ergeben sich allerdings noch keine signifikanten Vorteile im Betrieb der Anwendung. Um die mit IaaS erworbene Skalierbarkeit, Flexibilität und vereinfachte Verwaltung zu nutzen, bedarf es eines Neudenkens im Design der Infrastruktur und ihren Bereitstellungsprozessen, im allgemeinen als Cloud-Native-Architektur bezeichnet [Quelle]. Allerdings bringen diese neuen Konzepte wiederum Herausforderungen mit sich, die es zu identifizieren und bewältigen gilt. \newline 
Im Rahmen der Bachelorthesis werden bereits gängige Konzepte und \glqq Best-Practices\grqq{} im Betrieb von cloudnahen Anwendungen evaluiert. Am Beispiel der Startnext-Infrastruktur und durch das Erstellen von \textit{Proof of Concepts} werden sie auf ihre Anwendbarkeit, Sinnhaftigkeit auf ihre Machbarkeit untersucht und diskutiert. %willst du hier noch sagen, dass sie wenn möglich nach Umsetztung auch noch bewertet werden, so dass man sehen kann ob sie die erhofftne Erfolge bringen. 

%Ich würde noch mal den Bezug zu Startnext herstellen, damit klar wird, dass wir alle diese Maßnahmen bei SN durchführen als Fallbeispiel.

\section{\label{sec:fachliche_ueberlegungen} Fachliche Überlegungen und Zielsetzung}
\subsection{Fachliche Überlegungen}
Cloud-Native dient als Überbegriff für eine Vielzahl an Paradigmen und Prinzipien, die das Entwerfen und Implementieren von Infrastrukturen und Software als skalierbare, resiliente, administrierbare, observierbare und automatisierte Systeme ermöglichen. Eine übergeordnete Rolle in Cloud Native Architekturen spielt das \textit{Immutable Infrastructure} (ImmuInfra) Paradigma, welches die Grundlage für die Implementierung vieler weiterer Prinzipien bildet. Dieser Ansatz zur Software-Verwaltung verfolgt das Ziel eines weitestgehend unveränderlichen Zustands eines bestehenden Systems. Wird also ein Update eines Servers benötigt, wird ein neuer Server mit dem neuen gewünschten Zustand erstellt und die alte Instanz entfernt, anstatt Komponenten im bestehenden System einfach zu verändern \cite{Hirschfeld:2018}. 
Dieser Ansatz ist eine direkte Antwort auf das auseinanderdriften von Zuständen und dem dadurch unkalkulierbaren Ausgang von Zustandsänderungen wie in \cref{sec:problem_loesung} beschrieben. Zusätzlich müssen Änderungen nicht mehr an bereits produktiven Umgebungen vorgenommen werden, deren Ausfall direkte Auswirkungen auf die Anwendung hätte. Im Gegenteil können diese Änderungen ausgeführt, geprüft und anschließend in den produktiven Einsatz geschaltet werden.
\subsubsection{Technologien} 
Bei der Umsetzung von Immutable Infrastructure spielen Werkzeuge und Technologien eine entscheidende Rolle, um die gewünschte Nachvollziehbarkeit, Berechenbarkeit sowie Skalierbarkeit zu gewährleisten. \newline 
Um Systemumgebungen reproduzierbar zu erstellen, können Images genutzt werden. Ein Image ist Abbild des Speichers einer Maschine mit allen Metadaten und Konfigurationen, die benötigt werden, um die Maschine wiederherzustellen oder zu duplizieren. %Image find ich hier schwach definiert und ist Maschine überhaupt schon definiert oder erklärt worden?
Moderne Tools ermöglichen das Erstellen von Konfigurationen als Code (As-Code), aus denen Maschinenimages für unterschiedliche Plattformen erstellt werden können. Durch den \glqq As-Code\grqq{}-Ansatz werden Images versionierbar.%Das find ich auch missverständlich für As-Code. War Konfiguration jemals großartig unterschiedlich? Also vielleicht weil man jetzt tools hat die alle interactiven Sachen verskripten. Meinst du hier Packer? Meinst du provisionierungs tools wie Ansible? Infrastructure-As-Code wie terraform? Ich würde auch sagen, man behandelt config jetzt so wie source code. Ich glaube ich würde hier erst auf everything as code eingehen und die Vorteile von versionierung, Nachvollziehbarkeit und Automatisierung raus stellen und dann die Beispiele Packer, Ansible, Terraform durchgehen wie es schon fast machst.
Änderungen an diesen sind nachvollziehbar und widerruflich. Zusätzlich können über verschiedene Instanzen genutzte Grund-Images, auch Golden Images genannt \cite{GoldenImage:2018}, für eine einheitliche und erleichterte Wartung der Systeme sorgen. 
\newline
Der \textit{As-Code-Ansatz } ist ebenso bei der Provisionierung von Hosts möglich. Infrastructure-As-Code (IaC) Tools ermöglichen das Erstellen von Umgebungskonfigurationen als Code, die von IaaS-Anbietern interpretierbar sind. So wird Idempotenz bei der Bereitstellung von Hosts gewährleistet. Die Provisionierung wird sowohl nachvollziehbar als auch automatisierbar und Änderungen an Umgebungen sind versionier- sowie testbar.
Durch Versionierung und in Form von Code vorliegende Konfigurationen aller Bestandteile der Anwendung ist es möglich, den gesamten Bereitstellungsprozess zu automatisieren und den in \cref{sec:problem_loesung} beschriebenen Anforderungen an moderne Infrastruktur zu entsprechen. Der Prozess vom Provisionieren des Hosts, über das Konfigurieren dieser und dem Aufspielen der Software soll gemeinsam mit entsprechenden Tests jeder dieser Schritte automatisiert ablaufen. So kann auf Anforderungsänderungen schnell reagiert werden und komplexe Infrastrukturen bleiben administrierbar.

Darüber hinaus werden Konzepte und Tools benötigt, welche den Anwendungsbetrieb und ihre automatisierte Bereitstellung sicher, einfach zu administrieren und unanfällig für Fehler machen. Das Verwalten von Geheimnissen, Orchestrieren und Registrieren von Services und Testens sind hier zu nennen.

\section{Verwandte Arbeiten}

Die ansteigende Nachfrage nach Cloud-Native Technologien spiegelt sich auch in der wissenschaftlichen Literatur wieder. So findet das Thema Immutable Infrastructure als eine der grundlegenden Disziplinen in Cloud-Native hohe Aufmerksamkeit in wissenschaftlichen Abhandlungen.

Rob Hirschfeld diskutiert auf der SREcon-Konferenz 2018 Immutable Deployments und ihren Mehrwert \cite{Hirschfeld:2018}. Er arbeitet heraus, dass der Ansatz Immutable Infrastructure zusammen mit Cloud Computing zum kosteneffizienten Ansatz für das Bereitstellen von Laufzeitumgebungen wird. Hirschfeld teilt Anwendungsweisen von ImmuInfra in drei verschiedene Immutable Patterns ein:  
\begin{itemize}
\item Einfaches Zurücksetzen und Neubespielen der Maschine
\item Live-Booten der Maschine und anschließendes Konfigurieren
\item Erstellen eines fertigen Images und Aufspielen des Images
\end{itemize}

Diese drei Immutable Patterns gehen alle mit eigenen Vor- sowie Nachteilen einher.
Während der erste Ansatz keine vollständige ImmuInfra-Lösung ist, bietet er dennoch mehr Nachvollziehbarkeit als klassiche Deployments. Pattern 2 und 3 unterscheiden sich hauptsächlich in der Nutzung von Images, die eine Wiederverwendbarkeit, Zurückrollen und ein zeitlich effektives Aufspielen ermöglichen. Der Beitrag stellt klar, dass die Nutzung von Images die meisten Vorteile bietet, allerdings auch die höchste Komplexität mit sich bringt.

Mit der Gestaltung eines resilienten und skalierenden Cloud-Native Systems durch eine selbstadministrierende Anwendung beschäftigen sich Toffetti et al. in \cite{Toffetti:2017}. Die Autoren erstellen ein Cloud-Native Design, welches in der Lage ist, selbständig zu skalieren und sich selbst zu heilen. Die Architektur beinhaltet neben Cloud Orchestrierung und verteiltem Konfigurationsmanagement ein verteiltes Datenbanksystem sowie die Fähigkeit der Leader-Election. Anschließend wird dieses Design auf eine bestehende Legacy Web-App angewandt. 
Er und sein Team emulieren Fehler sowohl in den containerisierten Teilen der Anwendung als auch in den VMs, welche die Container hosten. Sie messen die Reaktion des Systems auf Fehler und diskutieren diese. Der Beitrag bietet einen robusten Leitfaden für die Wahl 
von DesignPatterns und eine Orientierung für das Erstellen eines selbstgemanagten Cloud-Native Systems.

%Eins, zwei diskutierende Sätze zum Journal.
\pagebreak

%---------------------------------
% KONKRETE AUFGABEN
%---------------------------------
\section{Konkrete Aufgaben und Zeitplan}

Die Ausführung der Bachelorarbeit ist zeitlich wie folgt geplant:
\newline
Dauer: 9 Wochen

\begin{itemize}
	\item Woche 1: Literaturrecherche, Einleitung, Einrichten des Dokuments
 	\item Woche 2: Verfassen der Problemstellung, Zielsetzung und möglichen Lösungsansätzen
 	\item Woche 3: Entwickeln der Anforderungsanalyse
 	\item Woche 4: Grundlagen und verwandte Arbeiten - Verfassen und Recherche
 	\item Woche 5: Konzeption - Entwickeln eines Entwurfs zur Lösung der Problemstellung
 	\item Woche 6-7: Prototypische Umsetzung
 	\item Woche 8: Evaluation - Bewertung und Diskussion der Ergebnisse
 	\item Woche 9: Überarbeitung, Korrektur und Druck des Dokuments
\end{itemize}


%---------------------------------
% VORLÄUFIGE GLIEDERUNG
%---------------------------------
\section{\label{sec:gliederung}Vorläufige Gliederung}

{\parindent=5mm 
Eigenständigkeitserklärung

Zusammenfassung / Abstract}

\begin{gliederung}
	\item Einleitung
	\begin{gliederung}
		\item Problemfeld
		\item Ziesetzung
		\item Lösungsansatz
		\item Aufbau der Arbeit
	\end{gliederung}

	\item Anforderungsanalyse
	\begin{gliederung}
		\item Funktionale Anforderungen
		\item Nicht-funktionale Anforderungen
		\item Zusammenfassung
	\end{gliederung}
	
	\item Grundlagen und verwandte Arbeiten
	\begin{gliederung}
		\item Immutable Infrastructure
			\item HashiCorp - Packer
		\item Selfmanaged Infrastructure
		\item Monitoring von Infrastruktur
		\item Automatisierte Orchestrierung und Konfiguration von Infrastruktur
			\item HashiCorp - Nomad \& Consul
		\item Verwandte Arbeiten
	\end{gliederung}
	
	\item Konzeption
	\begin{gliederung}
		\item Entwurf einer Cloud-Native Infrastruktur
		\begin{gliederung}
			\item Image-Erstellung
			\item Deployment
			\item Konfiguration
			\item Orchestrierung
			\item Monitoring
		\end{gliederung}
		\item Zusammenfassung
	\end{gliederung}
	
	\item Prototypische Realisierung
	\begin{gliederung}
		\item Grundlegender Aufbau der Infrastruktur
		\item Konfigurieren der Services
		\item Qualitätssicherung
		\item Zusammenfassung
	\end{gliederung}

	\item Evaluation
	\begin{gliederung}
		\item Überprüfung funktionaler Anforderungen
		\item Überprüfung nicht-funktionaler Anforderungen
		\item Zusammenfassung
	\end{gliederung}

	\item Zusammenfassung und Ausblick
	\begin{gliederung}
		\item Zusammenfassung
		\item Ausblick
	\end{gliederung}

\end{gliederung}

{\parindent=5mm
Literaturverzeichnis}%\\[0.25cm]
% Vertikaler Abstand
%\vspace*{\baselineskip}
%\smallskip  % kleiner variabler Abstand
\medskip    % mittlerer variabler Abstand
%\bigskip    % großer variabler Abstand

\pagebreak

%---------------------------------
% PERSONEN
%---------------------------------
\section{Unterschriften}

Student:

 Name: \textit{Lukas Karsten}
 
 Adresse: \textit{Bismarckstraße 156, 28205 Bremen}
 
 Telefon: \textit{0157/81733962}
  
 E -Mail: \textit{lkarsten@stud.hs-bremen.de}


 \hfill \rule{7cm}{0.2mm}

 \hfill Unterschrift (Student\_in)
 
\noindent
Erstgutachter\_in / Betreuer\_in: \textit{ Prof. Dr. Lars Braubach?/Hans Mündelein? } \\[0.2cm]

 \hfill \rule{7cm}{0.2mm}

 \hfill Unterschrift (Erstgutachter)


\noindent
Zweitgutachter\_in: \textit{}




%---------------------------------
% LITERATUR
%---------------------------------
% \bibliographystyle{IEEEtran}
%\bibliography{IEEEabrv,literatur}
%\bibliographystyle{plain}
% \bibliography{literatur}


% Gesamte Literaturliste
%\printbibliography	

% Literaturliste getrennt in Papier- und Online-Quellen:
\printbibheading
\printbibliography[nottype=online, heading=subbibliography, title={Gedruckte Quellen}]
\printbibliography[type=online, heading=subbibliography, title={Online Quellen}]


\end{document}
